\documentclass[11pt,reqno]{amsart}

%==============================================================
% First-Chaos Dual Closure of the IBM Ito Correction
% Sequel to: Energy-First Stochastic Calculus (EnergyFirst.tex)
% Scope: Tier-B only. Close the L^2-divergent first scenery-chaos
% sector of the mollified IBM correction in H^{-1}.
%==============================================================

\usepackage[utf8]{inputenc}
\usepackage[T1]{fontenc}
\usepackage{lmodern}
\usepackage[margin=1.15in]{geometry}
\usepackage{amsmath,amssymb,amsthm,mathtools}
\usepackage{mathrsfs}
\usepackage{microtype}
\usepackage{hyperref}
\usepackage{cleveref}

\emergencystretch=3em

\theoremstyle{plain}
\newtheorem{theorem}{Theorem}[section]
\newtheorem{proposition}[theorem]{Proposition}
\newtheorem{lemma}[theorem]{Lemma}
\newtheorem{corollary}[theorem]{Corollary}

\theoremstyle{definition}
\newtheorem{definition}[theorem]{Definition}
\newtheorem{remark}[theorem]{Remark}
\newtheorem{assumption}[theorem]{Assumption}

\theoremstyle{remark}
\newtheorem*{sketch}{Proof sketch}
\newtheorem*{risk}{Risk}

\newcommand{\R}{\mathbb{R}}
\newcommand{\N}{\mathbb{N}}
\newcommand{\E}{\mathbb{E}}
\newcommand{\Prob}{\mathbb{P}}
\newcommand{\Hilb}{\mathcal{H}}
\newcommand{\Dom}{\mathrm{Dom}}
\newcommand{\Energy}{\mathcal{E}}
\newcommand{\Source}{\mathsf{S}}
\newcommand{\Obs}{\mathsf{Z}}
\newcommand{\IBM}{Z}
\DeclareMathOperator{\Var}{Var}
\DeclareMathOperator{\Cov}{Cov}
\DeclareMathOperator{\supp}{supp}
\DeclareMathOperator{\sgn}{sgn}

\title{First-Chaos Dual Closure of the\\
       Iterated Brownian Motion It\^o Correction}
\author{Ramiro Fontes}
\date{\today}

\begin{document}

\begin{abstract}
For iterated Brownian motion $\IBM_t=X_{|Y_t|}$ with independent
two-sided Brownian scenery $X$ and Brownian clock $Y$, the companion
paper~\cite{Fontes-EnergyFirst} proved a mollified pathwise It\^o
identity together with the analytic obstruction that the first
scenery-chaos projection of the mollified It\^o correction
$C_{\IBM,\varepsilon}(f;t)$ has variance diverging to $+\infty$ in
$L^{2}(\Omega)$ as $\varepsilon\to 0$, driven by the rank-one
$\partial_{r}^{2}X^{\varepsilon}$ kernel and by the Bouleau--Yor
spatial quadratic variation of the Brownian local time.

We close that divergent sector in the natural \emph{negative Sobolev
/ quartic observable-energy dual} topology. Writing
$\mathrm{Proj}_{1}^{X}C_{\IBM,\varepsilon}(f;t)=I_{1}^{X}(K_{\varepsilon}(\cdot\mid Y))$
with the exact kernel decomposition
\[
 K_{\varepsilon}\;=\;K_{\varepsilon}^{\mathrm{bdy}}
 \;+\;\tfrac12\bigl(T_{1}^{\varepsilon}+T_{2}^{\varepsilon}+T_{3}^{\varepsilon}\bigr)
 \;+\;R_{\varepsilon}^{\mathrm{layer}},
\]
where the first three summands extend the
bulk-Stein-extrapolated terms of the companion paper to all of
$[0,\infty)$ and the residual $R_{\varepsilon}^{\mathrm{layer}}$
captures the $r\in[0,c\varepsilon]$ boundary layer where the Stein
cancellation fails, we prove:
\begin{enumerate}
\item (Interior, \cref{thm:interior}) For every $f\in\mathscr S(\R)$,
both the bulk-extrapolated kernel
$K_{\varepsilon}^{\mathrm{int}}\coloneqq\tfrac12(T_{1}^{\varepsilon}+T_{2}^{\varepsilon}+T_{3}^{\varepsilon})$
and the actual interior kernel
$K_{\varepsilon}^{\mathrm{int,exact}}\coloneqq K_{\varepsilon}^{\mathrm{int}}+R_{\varepsilon}^{\mathrm{layer}}$
converge to the same limit
\[
 K^{\mathrm{int}}(w\mid Y)=\tfrac12\bigl[T_{1}^{*}(w\mid Y)
 +B(w)L(w)-\partial_{w}\!\bigl(C(w)L(w)\bigr)\bigr]
\]
in $L^{2}(\Omega_{Y};H^{-1}(\R,dw))$, with the boundary-layer
residual vanishing in the stronger space
$L^{2}(\Omega_{Y};L^{2}(\R))$ at rate $O(\varepsilon^{1/2})$
(\cref{lem:bdy-layer-vanish}). Here $L=L_{t}^{\cdot}(|Y|)$, and
$A,B,C$ are the heat-semigroup coefficients
$\tfrac14P_{|\cdot|}f^{(5)}$, $P_{|\cdot|}f'''$, $P_{|\cdot|}f'$
from a self-contained Stein/covariance audit.
\item (Boundary, \cref{thm:boundary})
$K_{\varepsilon}^{\mathrm{bdy}}\to L_{t}^{0}(Y)\,f'(x)\,\delta_{0}$ in
$L^{2}(\Omega_{Y};H^{-1/2-\delta}(\R))$ at rate
$\varepsilon^{\delta\wedge 1/2}$ for every $\delta>0$. The
$\delta_{0}'$ contribution conjectured from naive Malliavin
expansion vanishes identically because
$\Cov(X_{0}^{\varepsilon},\partial_{r}X_{r}^{\varepsilon}|_{r=0})=0$
(\cref{lem:bdy-cov-zero}).
\item (Combined, \cref{cor:combined})
$K_{\varepsilon}\to K^{\mathrm{int}}+L_{t}^{0}(Y)f'(x)\delta_{0}$ in
$L^{2}(\Omega_{Y};H^{-1}(\R))$, using the continuous embedding
$H^{-1/2-\delta}\subset H^{-1}$ for $0<\delta<1/2$.
\item (Quantitative rate, \cref{prop:interior-rate}) Under
\cref{assu:lt-moment} below, the interior rate is
$\varepsilon^{1/2-\eta}$ for every $\eta\in(0,1/2)$.
\end{enumerate}

Under the EnergyFirst nondegeneracy hypothesis $(P_{\eta}f')(x)\not\equiv 0$
on $(0,\infty)$, this convergent $H^{-1}$ sector is precisely the
sector whose $L^{2}(\Omega)$-variance diverges, so the dual closure
above is sharp at the topology selected by the proved obstruction.

The label \emph{quartic} refers throughout to
operator-derivative / test-energy order, not Wiener chaos degree
(\cref{rem:quartic-not-chaos}); the singular object is genuinely
first-$X$-chaos, but it is a contraction channel of the quartic IBM
correction. We do not construct $C_{\IBM}(f;t)$ in the full
$\Energy_{\IBM}^{(4),-}$ here; the $(m,n)$-bidegrees with $m\ge 2$,
the full Wiener--It\^o assembly, and the semigroup correction
$u(t,x)-f(x)$ remain the Tier-A problem (\cref{sec:scope}).
Filtration-level obstructions (Tsirelson-type noises etc.) are an
independent direction and play no role here.
\end{abstract}

\maketitle

%==============================================================
\section{Introduction}\label{sec:intro}

\subsection*{The proved obstruction.}
For mollified IBM $\IBM_{t}^{\varepsilon}=x+X^{\varepsilon}_{|Y_{t}|}$
with $X^{\varepsilon}=\rho_{\varepsilon}*X$,
$f\in\mathscr S(\R)$, and
$g_{\varepsilon}(r)=f(x+X^{\varepsilon}_{r})$, the mollified pathwise
It\^o theorem of~\cite{Fontes-EnergyFirst} establishes the identity
(writing $\delta_{Y}$ for the Skorohod divergence on the clock
factor):
\begin{equation}\label{eq:flagship}
 f(\IBM_{t}^{\varepsilon})-f(\IBM_{0}^{\varepsilon})
  =\delta_{Y}\bigl(B_{f,\varepsilon}(t)\bigr)+C_{\IBM,\varepsilon}(f;t),
\end{equation}
\begin{equation}\label{eq:Ceps}
 C_{\IBM,\varepsilon}(f;t)
  =g_{\varepsilon}'(0)\,L_{t}^{0}(Y)
  +\tfrac12\!\int_{0}^{t}g_{\varepsilon}''(|Y_{s}|)\,ds,
\end{equation}
with $g_{\varepsilon}'(r)=f'(x+U_{r})V_{r}$ and
$g_{\varepsilon}''(r)=f''(x+U_{r})V_{r}^{2}+f'(x+U_{r})W_{r}$ for
$U_{r}=X^{\varepsilon}_{r}$, $V_{r}=\partial_{r}X^{\varepsilon}_{r}$,
$W_{r}=\partial_{r}^{2}X^{\varepsilon}_{r}$. The same paper proves
the obstruction
\begin{equation}\label{eq:obstruction}
 \liminf_{\varepsilon\to 0}\,
 \Var\bigl(\mathrm{Proj}_{1}^{X}\,C_{\IBM,\varepsilon}(f;t)\bigr)
  =+\infty
\end{equation}
under the nondegeneracy hypothesis
\begin{equation}\label{eq:nondegen}
 C(\eta)\coloneqq(P_{\eta}f')(x)\not\equiv 0
 \quad\text{on }(0,\infty),
\end{equation}
so the first-$X$-chaos projection of $C_{\IBM,\varepsilon}$ has no
$L^{2}(\Omega)$ limit. \Cref{eq:obstruction} is sharp at the
$L^{2}(\Omega)$ level, but it does not preclude a limit in a weaker,
\emph{structurally identified}, topology.

\subsection*{Strategy: where the dropped derivative goes.}
The proof of~\eqref{eq:obstruction} in~\cite{Fontes-EnergyFirst}
isolates a single offending piece:
$T_{3}^{\varepsilon}(w)=-\partial_{w}(\rho_{\varepsilon}*C_{\varepsilon}L)(w)$
with $C_{\varepsilon}\to C$ uniformly on compacts of $(0,\infty)$ and
$L=L_{t}^{\cdot}(|Y|)$. The $L^{2}(\R)$ blow-up is exactly
``$\chi C L\notin H^{1}(\R)$'', a Bouleau--Yor failure. Dropping one
spatial derivative absorbs the blow-up:
$L\in H^{1/2-\eta}_{\mathrm{loc}}$ a.s., so
$\partial_{w}(CL)\in H^{-1/2-\eta}\subset H^{-1}(\R)$ a.s.\ and
$T_{3}^{\varepsilon}\to-\partial_{w}(CL)$ in $H^{-1}$.

The same one-derivative drop on the spatial scenery variable absorbs
all three interior pieces
$T_{1}^{\varepsilon},T_{2}^{\varepsilon},T_{3}^{\varepsilon}$ and,
after a separate treatment in $H^{-1/2-\delta}$, the boundary piece
$K_{\varepsilon}^{\mathrm{bdy}}$.

\subsection*{Conceptual rule: ``quartic'' means operator order, not chaos degree.}
The companion paper labels the natural unmollified target
$C_{\IBM}(f;t)$ as living in the dual of a \emph{quartic composite
observable energy}
$\Energy_{\IBM}^{(4)}=\Energy_{\IBM}^{(2)}\otimes\Energy_{\IBM}^{(2)}$.
The label \emph{quartic} refers throughout to the order of the
operator derivative producing the kernel: the second derivative
$D_{\Source}^{(2)}$ acts twice in the assembly of $g''_{\varepsilon}$,
and the test side requires four-fold structure when paired with the
correction. The singular object \eqref{eq:obstruction} is genuinely
first-$X$-chaos in the Wiener decomposition, and the closure topology
we identify is the negative Sobolev space $H^{-1}(\R)$ on the
spatial scenery variable; this is a single \emph{contraction channel}
of the four-fold composite, not a four-variable kernel space. We
make this distinction definitional in~\cref{sec:test-dual} and refer
to it throughout (\cref{rem:quartic-not-chaos}).

\subsection*{Main results.}
The following are stated informally; precise hypotheses are in
\cref{sec:test-dual,sec:audit,sec:interior,sec:boundary}.
\begin{itemize}
\item (Interior, \cref{thm:interior}) For every $f\in\mathscr S(\R)$,
the interior portion of $\mathrm{Proj}_{1}^{X}\,C_{\IBM,\varepsilon}$
has a Malliavin kernel $K_{\varepsilon}^{\mathrm{int}}$ converging to
an explicit limit in $L^{2}(\Omega_{Y};H^{-1}(\R,dw))$.
\item (Quantitative rate, \cref{prop:interior-rate}) Under explicit
assumptions on the boundary-layer kernel and on local-time Sobolev
moments, the interior rate is $\varepsilon^{1/2-\eta}$ for every
$\eta\in(0,1/2)$.
\item (Boundary, \cref{thm:boundary}) The boundary portion converges
to $L_{t}^{0}(Y)\,f'(x)\,\delta_{0}$ in
$L^{2}\bigl(\Omega_{Y};\,H^{-1/2-\delta}(\R)\bigr)$ at rate
$\varepsilon^{\delta\wedge 1/2}$.
\item (Combined, \cref{cor:combined}) The full first-$X$-chaos
Malliavin kernel converges in $L^{2}(\Omega_{Y};H^{-1}(\R))$ via
the embedding $H^{-1/2-\delta}\subset H^{-1}$,
$\delta\in(0,1/2)$.
\item (Sharpness vs.\ obstruction) Under~\eqref{eq:nondegen}, the
$H^{-1}$-convergent first-$X$-chaos sector above is precisely the
sector whose $L^{2}(\Omega)$-variance diverges in
\cite{Fontes-EnergyFirst}.
\end{itemize}

\subsection*{What this paper does not prove.}\label{sec:not}
We do \emph{not} construct $C_{\IBM}(f;t)$ as an element of the full
quartic-energy dual $\Energy_{\IBM}^{(4),-}$. We do \emph{not}
recover the Brownian-time semigroup correction $u(t,x)-f(x)$ from
the first-chaos sector, because first chaos has mean zero; the
semigroup correction lives in the $(0,n)$ pure-clock and
$(2,n)$ second-energy-contraction sectors (in Gaussian coordinates,
the latter is the Wick-square sector), not addressed here. We do \emph{not} pursue
filtration-level obstructions (Tsirelson-type noises etc.); they are
a separate direction and are not used in any argument below. The
sole goal is to close the $L^{2}(\Omega)$-divergent first-scenery-chaos
sector identified by~\eqref{eq:obstruction} in the dual topology
that the structure of the obstruction selects.

\subsection*{Outline.}
\Cref{sec:test-dual} defines the first-chaos test/dual block.
\Cref{sec:audit} re-derives the Stein/covariance audit of the
mollified kernel from first principles, including the boundary
covariance identity $\Cov(U_{0},V_{0})=0$ that determines the
boundary limit. \Cref{sec:interior} proves the interior convergence
theorem (unconditional) and the conditional rate proposition.
\Cref{sec:boundary} proves the boundary theorem.
\Cref{sec:combined} states the combined first-chaos dual closure and
its scope. \Cref{sec:imports} records the lemmas imported from
\cite{Fontes-EnergyFirst}; \cref{sec:risks} lists the points that
must be verified.

%==============================================================
\section{The first-chaos test/dual block}\label{sec:test-dual}

We do not build the full quartic dual $\Energy_{\IBM}^{(4),-}$ here.
Only its first-$X$-chaos slice is needed.

\begin{definition}[First-chaos test space and dual]\label{def:first-chaos-block}
Set
\[
 \Energy_{\IBM,1,*}^{(4),+}\;\coloneqq\;H^{1}(\R,dw),
 \qquad
 \Energy_{\IBM,1,*}^{(4),-}\;\coloneqq\;H^{-1}(\R,dw),
\]
the standard Bessel-potential Hilbert spaces on the scenery variable
$w\in\R$, with duality pairing
$\langle K,\phi\rangle_{H^{-1},H^{1}}$ extending the
$L^{2}(\R)$-inner product. The \emph{first-chaos pairing} between a
first-$X$-chaos random variable $F=I_{1}^{X}(K)$ with conditional
Malliavin kernel $K(\cdot\mid Y)\in L^{2}(\Omega_{Y};H^{-1}(\R))$
and a deterministic test integrand $\phi\in H^{1}(\R)$ is
\[
 \bigl\langle F,\phi\bigr\rangle_{1,*}
 \;\coloneqq\;\E_{Y}\!\bigl[\langle K(\cdot\mid Y),\phi\rangle_{H^{-1},H^{1}}\bigr],
\]
finite by Cauchy--Schwarz once $K(\cdot\mid Y)\in H^{-1}(\R)$ a.s.\ in
$L^{2}(\Omega_{Y})$. When $K(\cdot\mid Y)\in L^{2}(\R)$ a.s.\ the
duality pairing reduces to the $L^{2}(\R)$-pairing.
\end{definition}

\begin{remark}[``Quartic'' = operator order, not Wiener-chaos degree]
\label{rem:quartic-not-chaos}
\Cref{def:first-chaos-block} concerns elements of the first scenery
Wiener chaos $\Hilb_{1}^{X}$. The label \emph{quartic} attached to
the dual $\Energy_{\IBM,1,*}^{(4),-}$ refers to the order of the
operator-derivative composition that produced the kernel
($D_{\Source}^{(2)}$ acting twice in the assembly of
$g_{\varepsilon}''$, paired against second-order tensor test
integrands), \emph{not} the Wiener chaos degree of the resulting
random variable. The first-scenery-chaos sector is a single
contraction channel of the four-fold composite identified in the
rigged-quartic hierarchy of~\cite{Fontes-EnergyFirst}, and is the
only sector closed in this paper.
\end{remark}

\begin{remark}[Energy contraction vs.\ Wick]\label{rem:energy-contraction}
In the Gaussian scenery used for IBM, the contractions appearing
below may be represented by Wiener-chaos/Wick bookkeeping. However
the mechanism used here is weaker: it is the subtraction or
isolation of lower-order energy contractions determined by the
observable-energy pairing. Thus Wick ordering is not a primitive
of the framework; it is the Gaussian realization of the more
intrinsic energy-contraction operation. We accordingly use the
phrases \emph{energy contraction}, \emph{observable contraction},
\emph{compensated energy product}, \emph{lower-order contraction
subtraction}, and \emph{contraction channel} throughout, reserving
the word \emph{Wick} for explicit references to classical Gaussian
Wiener-chaos bookkeeping (e.g.\ when comparing with the
contraction-renormalised correction candidate
of~\cite{Fontes-EnergyFirst}).
\end{remark}

\begin{remark}[Why \texorpdfstring{$H^{-1}$}{H\textasciicircum-1} is a robust energy-dual block]\label{rem:why-Hminus1}
$H^{-1}(\R)$ is a robust, canonical energy-dual block for the
first-chaos sector because two minimal facts are available:
$\partial_{w}\colon L^{2}(\R)\to H^{-1}(\R)$ is bounded, and
$\rho_{\varepsilon}*C_{\varepsilon}L\to CL$ in $L^{2}(\R)$ for the
offending kernel of~\eqref{eq:obstruction}
in~\cite{Fontes-EnergyFirst}. Together these close the singular
channel in $H^{-1}$. We do not claim $H^{-1}$ is the optimal Sobolev
target: the sharper local threshold is governed by
Bouleau--Yor regularity, which puts $L=L_{t}^{\cdot}(|Y|)\in H^{1/2-\eta}_{\mathrm{loc}}$
and hence $\partial_{w}(CL)\in H^{-1/2-\eta}_{\mathrm{loc}}\subset H^{-1}_{\mathrm{loc}}$
a.s. The point is structural: $L^{2}(\Omega)$ is too strong to host
the unmollified first-chaos correction, while the negative
Sobolev / energy-dual topology closes the singular channel.

In particular, ordinary Wick-square renormalisation
(the Gaussian-coordinate realisation of a second-energy-contraction
subtraction) does not remove the rank-one first-chaos singularity
carried by $\partial_{r}^{2}X^{\varepsilon}$; the correct repair is
negative Sobolev / quartic-dual closure.
\end{remark}

\begin{remark}[Clock variable; full quartic space]
We do not introduce a clock-side regularity index here. The
first-chaos closure is entirely about the scenery variable $w$; the
clock variable enters only through the random $L_{t}^{\cdot}(|Y|)$
coefficient, expectation in $Y$ replaces clock-Sobolev pairing. A
clock-graded test space becomes necessary only at the $(m,n)$
mixed-bidegree level, deferred to Tier-A.
\end{remark}

%==============================================================
\section{Stein and covariance audit; the kernel}\label{sec:audit}

We re-derive the first-chaos Malliavin kernel of
$\mathrm{Proj}_{1}^{X}\,C_{\IBM,\varepsilon}(f;t)$ from first
principles. No coefficient is trusted from~\cite{Fontes-EnergyFirst}
until re-checked.

\subsection*{Setup.}
$X$ is two-sided Brownian scenery on $\R$, white-noise representation
$X_{z}=\int_{\R}\sigma_{z}(w)\,d\xi(w)$ with
$\sigma_{z}(w)=\sgn(z)\mathbf 1_{[0,z]_{*}}(w)$ and
$\Cov(X_{z},X_{z'})=\tfrac12(|z|+|z'|-|z-z'|)$. The mollifier $\rho$
is even, non-negative, $C_{c}^{\infty}$ with support in $[-c,c]$,
$\int\rho=1$;
$\rho_{\varepsilon}(\cdot)=\varepsilon^{-1}\rho(\cdot/\varepsilon)$,
$c_{\rho}^{(1)}\coloneqq\int(\rho')^{2}$,
$c_{\rho}\coloneqq\int|v|\rho(v)\,dv$.

\subsection*{Mollified scenery: Malliavin kernels.}
$U_{r}=X^{\varepsilon}_{r}=\int\rho_{\varepsilon}(r-z)X_{z}\,dz$,
$V_{r}=\partial_{r}U_{r}$, $W_{r}=\partial_{r}^{2}U_{r}$.
By integration by parts in the white-noise representation
($\rho_{\varepsilon}$ compactly supported),
\begin{equation}\label{eq:VW-rep}
 V_{r}=\int\rho_{\varepsilon}(r-z)\,d\xi(z),
 \qquad
 W_{r}=\int\rho_{\varepsilon}'(r-z)\,d\xi(z),
\end{equation}
so the Malliavin derivatives are
\begin{equation}\label{eq:Dw}
 D_{w}U_{r}=\xi_{r,\varepsilon}(w),\quad
 D_{w}V_{r}=\rho_{\varepsilon}(r-w),\quad
 D_{w}W_{r}=\rho_{\varepsilon}'(r-w),
\end{equation}
with
$\xi_{r,\varepsilon}(w)=\int\rho_{\varepsilon}(r-z)\sigma_{z}(w)\,dz$.

\subsection*{Bulk covariances.}
$\Var(U_{r})=\sigma_{\varepsilon}^{2}(r)$ with
$\sigma_{\varepsilon}^{2}(r)=
\int\rho_{\varepsilon}(r-z)|z|\,dz
-\tfrac12\iint\rho_{\varepsilon}(r-z_{1})\rho_{\varepsilon}(r-z_{2})|z_{1}-z_{2}|\,dz_{1}dz_{2}$.
For $|r|>c\varepsilon$, the first integral equals $|r|$ and the
second equals $\tfrac12 K_{\rho}\varepsilon$ for a finite constant
$K_{\rho}$ independent of $r$; hence
$\sigma_{\varepsilon}^{2}(r)=|r|+O(\varepsilon)$ uniformly in
$|r|>c\varepsilon$. Differentiating, in the bulk:
\begin{equation}\label{eq:bulk-cov}
 (\sigma_{\varepsilon}^{2})'(r)=1,\qquad
 \Cov(U_{r},V_{r})=\tfrac12,\qquad
 \sigma_{V}^{2}=c_{\rho}^{(1)}/\varepsilon,\qquad
 \Cov(U_{r},W_{r})=-c_{\rho}^{(1)}/\varepsilon.
\end{equation}
The last identity follows by differentiating
$2\Cov(U_{r},V_{r})=(\sigma_{\varepsilon}^{2})'(r)=1$ once more:
$2(\sigma_{V}^{2}+\Cov(U_{r},W_{r}))=0$.

\subsection*{Gaussian Stein identities (re-derived).}
For jointly centred Gaussian $(U,V)$ and smooth $h$,
$\E[h(U)V]=\Cov(U,V)\E[h'(U)]$ (Stein, direct). Iterating once on
$g(U,V)=h(U)V$ via the vector Stein identity
$\E[Vg(U,V)]=\Cov(V,U)\E[\partial_{U}g]+\Var(V)\E[\partial_{V}g]$
gives
\begin{equation}\label{eq:Stein-sq}
 \E[h(U)V^{2}]=\Cov(U,V)^{2}\E[h''(U)]+\sigma_{V}^{2}\E[h(U)].
\end{equation}

\subsection*{Bulk Stein cancellation.}
Apply $D_{w}$ to
$g_{\varepsilon}''(r)=f''(x+U_{r})V_{r}^{2}+f'(x+U_{r})W_{r}$
and take expectation in $X$:
\begin{align*}
 \E_{X}[D_{w}g_{\varepsilon}''(r)]
 &=\E_{X}[f'''(x+U_{r})V_{r}^{2}+f''(x+U_{r})W_{r}]\,\xi_{r,\varepsilon}(w)\\
 &\quad+2\E_{X}[f''(x+U_{r})V_{r}]\,\rho_{\varepsilon}(r-w)
 +\E_{X}[f'(x+U_{r})]\,\rho_{\varepsilon}'(r-w).
\end{align*}
In the bulk, \eqref{eq:Stein-sq} and~\eqref{eq:bulk-cov} give
$\E_{X}[f'''(x+U_{r})V_{r}^{2}]=\sigma_{V}^{2}\E_{X}[f'''(x+U_{r})]+\tfrac14\E_{X}[f^{(5)}(x+U_{r})]$
and
$\E_{X}[f''(x+U_{r})W_{r}]=\Cov(U_{r},W_{r})\E_{X}[f'''(x+U_{r})]
=-\sigma_{V}^{2}\E_{X}[f'''(x+U_{r})]$,
so the $\varepsilon^{-1}$ contributions cancel exactly. Setting
\begin{equation}\label{eq:ABC}
 A_{\varepsilon}(r)=\tfrac14\E_{X}[f^{(5)}(x+U_{r})],\quad
 B_{\varepsilon}(r)=\E_{X}[f'''(x+U_{r})],\quad
 C_{\varepsilon}(r)=\E_{X}[f'(x+U_{r})],
\end{equation}
we obtain, in the bulk,
\begin{equation}\label{eq:kernel-form}
 \E_{X}[D_{w}g_{\varepsilon}''(r)]
 =A_{\varepsilon}(r)\xi_{r,\varepsilon}(w)+B_{\varepsilon}(r)\rho_{\varepsilon}(r-w)
 +C_{\varepsilon}(r)\rho_{\varepsilon}'(r-w).
\end{equation}
Each $A_{\varepsilon},B_{\varepsilon},C_{\varepsilon}$ is uniformly
bounded in $(r,\varepsilon)$ with pointwise heat-semigroup limits
\begin{equation}\label{eq:ABC-limit}
 A(r)=\tfrac14(P_{|r|}f^{(5)})(x),\quad
 B(r)=(P_{|r|}f''')(x),\quad
 C(r)=(P_{|r|}f')(x),
\end{equation}
the limit $C(r)$ matching~\eqref{eq:nondegen}.

\subsection*{Boundary covariance:
\texorpdfstring{$\Cov(U_{0},V_{0})=0$}{Cov(U\_0,V\_0)=0}.}

\begin{lemma}[Boundary covariance vanishes]\label{lem:bdy-cov-zero}
For two-sided Brownian scenery $X$ and even mollifier $\rho$,
\[
 \Cov(U_{0},V_{0})\;=\;\tfrac12(\sigma_{\varepsilon}^{2})'(0)\;=\;0,
 \qquad\text{hence } U_{0}\perp V_{0}.
\]
\end{lemma}

\begin{proof}
$(\sigma_{\varepsilon}^{2})'(r)=I_{1}'(r)-\tfrac12 I_{2}'(r)$ with
$I_{1}(r)=\int\rho_{\varepsilon}(r-z)|z|\,dz$ and
$I_{2}(r)=\iint\rho_{\varepsilon}(r-z_{1})\rho_{\varepsilon}(r-z_{2})|z_{1}-z_{2}|\,dz_{1}dz_{2}$.
At $r=0$, by translation,
$I_{1}'(0)=-\int\rho_{\varepsilon}'(-z)|z|\,dz=\int\rho_{\varepsilon}'(z)|z|\,dz$,
which vanishes by odd-even pairing. For $I_{2}'(0)$, the integrand
is symmetric in $(z_{1},z_{2})$, so
$I_{2}'(0)=-2\iint\rho_{\varepsilon}'(-z_{1})\rho_{\varepsilon}(-z_{2})|z_{1}-z_{2}|\,dz_{1}dz_{2}
=2\iint\rho_{\varepsilon}'(z_{1})\rho_{\varepsilon}(z_{2})|z_{1}-z_{2}|\,dz_{1}dz_{2}$.
Integration by parts in $z_{1}$, using compact support of
$\rho_{\varepsilon}$, gives
$I_{2}'(0)=-2\iint\rho_{\varepsilon}(z_{1})\rho_{\varepsilon}(z_{2})\sgn(z_{1}-z_{2})\,dz_{1}dz_{2}=0$
by antisymmetry of $\sgn(z_{1}-z_{2})$ against the symmetric
product measure. Hence $(\sigma_{\varepsilon}^{2})'(0)=0$. Joint
Gaussianity of $(U_{0},V_{0})$ and vanishing covariance imply
independence.
\end{proof}

\begin{remark}[Consequence for the boundary kernel]
\label{rem:bdy-coeff}
Because $V_{0}\perp U_{0}$, $\E_{X}[f''(x+U_{0}^{\varepsilon})V_{0}]
=\E[V_{0}]\,\E_{X}[f''(x+U_{0}^{\varepsilon})]=0$. The
``$\xi_{0,\varepsilon}$ piece'' in $\E_{X}[D_{w}g_{\varepsilon}'(0)]$
vanishes identically, leaving only the $\rho_{\varepsilon}$ kernel
with $C_{\varepsilon}(0)$ coefficient. This rules out the naively
conjectured $\delta'_{0}$ term in the limit of
$K_{\varepsilon}^{\mathrm{bdy}}$; see \cref{thm:boundary}.
\end{remark}

\subsection*{Kernel decomposition.}
Combining the Stein audit with the IBM identity~\eqref{eq:Ceps} and
the occupation-density formula
$\int_{0}^{t}h(|Y_{s}|)\,ds=\int_{0}^{\infty}h(r)L_{t}^{r}(|Y|)\,dr$
($L\coloneqq L_{t}^{\cdot}(|Y|)$), the \emph{exact} conditional
Malliavin kernel of $\mathrm{Proj}_{1}^{X}C_{\IBM,\varepsilon}(f;t)$
in $L^{2}(\R,dw)$ is
\begin{equation}\label{eq:Keps-decomp}
 K_{\varepsilon}(w\mid Y)
 =K_{\varepsilon}^{\mathrm{bdy}}(w\mid Y)
 +\tfrac12\bigl(T_{1}^{\varepsilon}+T_{2}^{\varepsilon}+T_{3}^{\varepsilon}\bigr)(w\mid Y)
 +R_{\varepsilon}^{\mathrm{layer}}(w\mid Y),
\end{equation}
with
\begin{align}
 K_{\varepsilon}^{\mathrm{bdy}}(w\mid Y)
 &=L_{t}^{0}(Y)\cdot C_{\varepsilon}(0)\,\rho_{\varepsilon}(w),\label{eq:Kbdy}\\
 T_{1}^{\varepsilon}(w\mid Y)
 &=\int_{0}^{\infty}A_{\varepsilon}(r)L(r)\,\xi_{r,\varepsilon}(w)\,dr,\label{eq:T1}\\
 T_{2}^{\varepsilon}(w\mid Y)
 &=\bigl(\rho_{\varepsilon}*B_{\varepsilon}L\bigr)(w),\label{eq:T2}\\
 T_{3}^{\varepsilon}(w\mid Y)
 &=-\partial_{w}\bigl(\rho_{\varepsilon}*C_{\varepsilon}L\bigr)(w),\label{eq:T3}
\end{align}
where the $T_{j}^{\varepsilon}$ are defined using the
\emph{bulk-extrapolated} coefficients $A_{\varepsilon},B_{\varepsilon},
C_{\varepsilon}$ of~\eqref{eq:ABC} integrated over all of
$[0,\infty)$, and the \emph{boundary-layer residual} is
\begin{equation}\label{eq:R-layer}
 R_{\varepsilon}^{\mathrm{layer}}(w\mid Y)
 =\tfrac12\!\int_{0}^{c\varepsilon}\!\!\bigl[\Delta_{A}(r)\,\xi_{r,\varepsilon}(w)
 +\Delta_{B}(r)\,\rho_{\varepsilon}(r-w)\bigr]L(r)\,dr,
\end{equation}
with
\begin{equation}\label{eq:DeltaAB}
\begin{aligned}
 \Delta_{A}(r)&=\bigl(\sigma_{V}^{2}+\Cov(U_{r},W_{r})\bigr)\E_{X}[f'''(x+U_{r})]
 +\bigl(\Cov(U_{r},V_{r})^{2}-\tfrac14\bigr)\E_{X}[f^{(5)}(x+U_{r})],\\
 \Delta_{B}(r)&=\bigl(2\Cov(U_{r},V_{r})-1\bigr)\E_{X}[f'''(x+U_{r})],
\end{aligned}
\end{equation}
both supported on $r\in[0,c\varepsilon]$ (zero in the bulk by the
cancellation~\eqref{eq:kernel-form}).
(In~\eqref{eq:Kbdy} we have used
$\rho_{\varepsilon}(0-w)=\rho_{\varepsilon}(w)$ by evenness, and
\cref{lem:bdy-cov-zero} to kill the $\xi_{0,\varepsilon}$
contribution at $r=0$.)

\subsection*{Boundary layer of the bulk identity.}\label{ssec:boundary-layer}
The cancellation~\eqref{eq:kernel-form} that yields
$T_{1}^{\varepsilon},T_{2}^{\varepsilon},T_{3}^{\varepsilon}$ from
the bulk coefficients $A_{\varepsilon},B_{\varepsilon},C_{\varepsilon}$
uses $(\sigma_{\varepsilon}^{2})'(r)=1$, which holds only for
$|r|>c\varepsilon$. The residual $R_{\varepsilon}^{\mathrm{layer}}$
defined by~\eqref{eq:R-layer}--\eqref{eq:DeltaAB} captures the
deviation in the interior layer $r\in[0,c\varepsilon]$. Its
vanishing in $L^{2}(\Omega_{Y};L^{2}(\R))$, with the explicit
$\varepsilon^{1/2}$ rate~\eqref{eq:R-layer-bound} used in
\cref{prop:interior-rate}, is the content of
\cref{lem:bdy-layer-vanish} below.

%==============================================================
\section{Interior first-chaos convergence}\label{sec:interior}

We separate the qualitative convergence statement (unconditional,
\cref{thm:interior}, supported by the qualitative
\cref{lem:bdy-layer-vanish}) from the quantitative rate
(\cref{prop:interior-rate}, conditional on two stronger
assumptions).

\subsection*{Qualitative boundary-layer vanishing.}

\begin{lemma}[Boundary-layer residual vanishing]\label{lem:bdy-layer-vanish}
For every $f\in\mathscr S(\R)$, the boundary-layer residual
$R_{\varepsilon}^{\mathrm{layer}}$ of~\eqref{eq:R-layer}--\eqref{eq:DeltaAB}
satisfies
\[
 R_{\varepsilon}^{\mathrm{layer}}\;\xrightarrow[\varepsilon\to 0]{}\;0
 \quad\text{in }L^{2}\bigl(\Omega_{Y};\,L^{2}(\R,dw)\bigr),
\]
and hence in $L^{2}(\Omega_{Y};H^{-1}(\R))$ via the continuous
embedding $L^{2}\hookrightarrow H^{-1}$.
\end{lemma}

\begin{proof}
\emph{Step 1: closed form for $(\sigma_{\varepsilon}^{2})''$.}
Writing $\sigma_{\varepsilon}^{2}(r)=\E[(X_{r}^{\varepsilon})^{2}]$
and using $\Cov(X_{z_{1}},X_{z_{2}})=\tfrac12(|z_{1}|+|z_{2}|-|z_{1}-z_{2}|)$
for two-sided Brownian scenery, change variables
$u_{i}=r-z_{i}$ in
$\sigma_{\varepsilon}^{2}(r)=\iint\rho_{\varepsilon}(r-z_{1})\rho_{\varepsilon}(r-z_{2})\Cov(X_{z_{1}},X_{z_{2}})\,dz_{1}dz_{2}$
to obtain
\begin{equation}\label{eq:sigma2-closed-form}
 \sigma_{\varepsilon}^{2}(r)
 \;=\;\bigl(\rho_{\varepsilon}*|\cdot|\bigr)(r)\,-\,K_{\rho}\,\varepsilon,
 \qquad
 K_{\rho}\;\coloneqq\;\tfrac12\!\iint\rho(v_{1})\rho(v_{2})|v_{1}-v_{2}|\,dv_{1}dv_{2}.
\end{equation}
Differentiating the smooth convolution
$(\rho_{\varepsilon}*|\cdot|)$ twice in $r$ gives
$(\sigma_{\varepsilon}^{2})'(r)=(\rho_{\varepsilon}*\sgn)(r)$ and
\begin{equation}\label{eq:sigma2-second-derivative}
 (\sigma_{\varepsilon}^{2})''(r)\;=\;2\,\rho_{\varepsilon}(r),
\end{equation}
using $(\sgn)'=2\delta_{0}$ in the distributional sense (the
convolution with $\rho_{\varepsilon}\in C_{c}^{\infty}$ regularises
to a smooth function on $\R$).

\emph{Step 2: layer factor bounds.}
From~\eqref{eq:sigma2-second-derivative} and
$2(\sigma_{V}^{2}+\Cov(U_{r},W_{r}))=(\sigma_{\varepsilon}^{2})''(r)$,
\begin{equation}\label{eq:sigmaV-W-bound}
 \bigl|\sigma_{V}^{2}+\Cov(U_{r},W_{r})\bigr|
 \;=\;\rho_{\varepsilon}(r)
 \;\le\;\|\rho\|_{\infty}/\varepsilon
 \qquad\text{for }r\in[0,c\varepsilon],
\end{equation}
and the left-hand side vanishes for $r>c\varepsilon$, confirming
the bulk cancellation~\eqref{eq:kernel-form}. Similarly
$2\Cov(U_{r},V_{r})=(\sigma_{\varepsilon}^{2})'(r)=(\rho_{\varepsilon}*\sgn)(r)\in[-1,1]$,
and by oddness of $\sgn$ and evenness of $\rho_{\varepsilon}$ takes
value $0$ at $r=0$ (\cref{lem:bdy-cov-zero}); for $r\in[0,c\varepsilon]$
it transitions monotonically from $0$ to $1$, so
$\Cov(U_{r},V_{r})\in[0,1/2]$, hence
$|\Cov(U_{r},V_{r})^{2}-\tfrac14|\le\tfrac14$ and
$|2\Cov(U_{r},V_{r})-1|\le 1$. For $f\in\mathscr S(\R)$, every
derivative $f^{(k)}$ is bounded, and
$|\E_{X}[f^{(k)}(x+U_{r})]|\le\|f^{(k)}\|_{\infty}$. Combining,
for $r\in[0,c\varepsilon]$,
\begin{equation}\label{eq:Delta-bounds}
 |\Delta_{A}(r)|\;\le\;C_{f}\,\rho_{\varepsilon}(r)+\tfrac14\|f^{(5)}\|_{\infty},
 \qquad
 |\Delta_{B}(r)|\;\le\;\|f'''\|_{\infty},
\end{equation}
with $C_{f}=\|f'''\|_{\infty}$.

\emph{Step 3: $L^{2}(\R)$ bound on $R_{\varepsilon}^{\mathrm{layer}}$.}
The scenery-derivative kernel estimate
of~\cite{Fontes-EnergyFirst} (mollifier $L^{2}$-rate and kernel
bound) gives
$\|\xi_{r,\varepsilon}\|_{L^{2}(\R)}^{2}\le|r|+C_{\rho}\varepsilon$,
so
$\|\xi_{r,\varepsilon}\|_{L^{2}(\R)}\le(c+C_{\rho})^{1/2}\,\varepsilon^{1/2}$
for $r\in[0,c\varepsilon]$. Also
$\|\rho_{\varepsilon}\|_{L^{2}(\R)}=(\int\rho^{2})^{1/2}\,\varepsilon^{-1/2}$.
Apply the triangle inequality in $L^{2}(\R,dw)$ to the
representation~\eqref{eq:R-layer}, writing $L=L_{t}^{\cdot}(|Y|)$
and $\|L\|_{\infty}\coloneqq\sup_{r\ge 0}L_{t}^{r}(|Y|)$:
\begin{align*}
 \|R_{\varepsilon}^{\mathrm{layer}}\|_{L^{2}(\R)}
 &\le\tfrac12\!\int_{0}^{c\varepsilon}\!|\Delta_{A}(r)|L(r)\|\xi_{r,\varepsilon}\|_{L^{2}}\,dr
  +\tfrac12\!\int_{0}^{c\varepsilon}\!|\Delta_{B}(r)|L(r)\|\rho_{\varepsilon}\|_{L^{2}}\,dr\\
 &\le\tfrac12(c+C_{\rho})^{1/2}\varepsilon^{1/2}\|L\|_{\infty}\!\int_{0}^{c\varepsilon}\!|\Delta_{A}(r)|\,dr
  +\tfrac12\bigl(\textstyle\int\rho^{2}\bigr)^{1/2}\varepsilon^{-1/2}\|L\|_{\infty}\!\int_{0}^{c\varepsilon}\!|\Delta_{B}(r)|\,dr.
\end{align*}
Using~\eqref{eq:Delta-bounds} and
$\int_{0}^{c\varepsilon}\rho_{\varepsilon}(r)\,dr\le 1$,
$\int_{0}^{c\varepsilon}1\,dr=c\varepsilon$,
\[
 \int_{0}^{c\varepsilon}|\Delta_{A}(r)|\,dr\le C_{f}+\tfrac14\|f^{(5)}\|_{\infty}\,c\varepsilon,
 \qquad
 \int_{0}^{c\varepsilon}|\Delta_{B}(r)|\,dr\le\|f'''\|_{\infty}\,c\varepsilon.
\]
Substituting,
\begin{equation}\label{eq:R-layer-bound}
 \|R_{\varepsilon}^{\mathrm{layer}}\|_{L^{2}(\R)}
 \;\le\;C(f,\rho)\,\varepsilon^{1/2}\,\|L\|_{\infty}
 \qquad\text{a.s.\ in }Y,
\end{equation}
for some finite $C(f,\rho)$ depending only on
$\|f'''\|_{\infty},\|f^{(5)}\|_{\infty}$, on $c$, on $\|\rho\|_{\infty}$,
and on $\int\rho^{2}$.

\emph{Step 4: $L^{2}(\Omega_{Y})$ assembly and conclusion.}
Squaring~\eqref{eq:R-layer-bound} and taking $\E_{Y}$,
\[
 \E_{Y}\!\bigl[\,\|R_{\varepsilon}^{\mathrm{layer}}\|_{L^{2}(\R)}^{2}\,\bigr]
 \;\le\;C(f,\rho)^{2}\,\varepsilon\,\E_{Y}[\,\|L\|_{\infty}^{2}\,].
\]
The spatial supremum $\|L\|_{\infty}=\sup_{r\ge 0}L_{t}^{r}(|Y|)$
of the spatial Brownian local time at the reflected Brownian motion
$|Y|$ has finite moments of every order: writing
$L_{t}^{r}(|Y|)=L_{t}^{r}(Y)+L_{t}^{-r}(Y)$ for $r>0$ and using the
finite-moment estimate
$\E[\sup_{x\in\R}L_{t}^{x}(Y)^{p}]<\infty$ for one-sided Brownian
local time (\cite[Theorem 9.4.1]{MarcusRosen}), we obtain
$\E_{Y}[\|L\|_{\infty}^{2}]\le 4\,\E[\sup_{x\in\R}L_{t}^{x}(Y)^{2}]<\infty$.
Hence
$\E_{Y}\|R_{\varepsilon}^{\mathrm{layer}}\|_{L^{2}(\R)}^{2}\to 0$ at
rate $O(\varepsilon)$ as $\varepsilon\to 0$, proving convergence in
$L^{2}(\Omega_{Y};L^{2}(\R))$. The $L^{2}(\Omega_{Y};H^{-1}(\R))$
statement is immediate from the continuous embedding
$L^{2}(\R)\hookrightarrow H^{-1}(\R)$ with operator norm $\le 1$.
\end{proof}

\subsection*{Quantitative input (for the rate).}

The boundary-layer ingredient needed by the rate
\cref{prop:interior-rate} is already provided in explicit form by
\cref{lem:bdy-layer-vanish}: the proof
yields~\eqref{eq:R-layer-bound}, i.e.\
\[
 \|R_{\varepsilon}^{\mathrm{layer}}\|_{L^{2}(\R)}
 \;\le\;C(f,\rho)\,\varepsilon^{1/2}\,\|L_{t}^{\cdot}(|Y|)\|_{\infty}
 \quad\text{a.s.\ in }Y,
\]
with $C(f,\rho)$ depending only on $\|f'''\|_{\infty},\|f^{(5)}\|_{\infty}$,
on the support constant $c$, and on $\|\rho\|_{\infty},\int\rho^{2}$.
Together with $\E_{Y}[\|L_{t}^{\cdot}(|Y|)\|_{\infty}^{2}]<\infty$
(Marcus--Rosen, used in Step~4 of the proof of
\cref{lem:bdy-layer-vanish}) this gives an
$L^{2}(\Omega_{Y};L^{2}(\R))$-bound of size $O(\varepsilon^{1/2})$
on $R_{\varepsilon}^{\mathrm{layer}}$, with no further assumption.

The only quantitative input not already proved here is the spatial
Sobolev regularity of the local time:

\begin{assumption}[Local-time Sobolev moments]\label{assu:lt-moment}
For every $\eta\in(0,1/2)$ and $p\ge 2$,
\[
 \E_{Y}\bigl[\,\|L_{t}^{\cdot}(|Y|)\|_{H^{1/2-\eta}([0,M_{t}])}^{p}\,\bigr]<\infty.
\]
\end{assumption}

\Cref{assu:lt-moment} is a Marcus--Rosen-type moment estimate for
the spatial Sobolev regularity of the reflected-BM local time
(\cite{MarcusRosen}, Theorem 9.4.1 and Chapter 10). Qualitative
convergence (\cref{thm:interior}) needs only the proved
\cref{lem:bdy-layer-vanish}; the $\varepsilon^{1/2-\eta}$ rate
(\cref{prop:interior-rate}) additionally needs
\cref{assu:lt-moment}.

\subsection*{Qualitative convergence.}

Define the \emph{bulk-extrapolated interior kernel} and the
\emph{actual interior kernel}, respectively, as
\[
 K_{\varepsilon}^{\mathrm{int}}\;\coloneqq\;\tfrac12\bigl(T_{1}^{\varepsilon}+T_{2}^{\varepsilon}+T_{3}^{\varepsilon}\bigr),
 \qquad
 K_{\varepsilon}^{\mathrm{int,exact}}\;\coloneqq\;K_{\varepsilon}^{\mathrm{int}}+R_{\varepsilon}^{\mathrm{layer}};
\]
$K_{\varepsilon}^{\mathrm{int,exact}}$ is the actual integrand-derived
contribution to~\eqref{eq:Keps-decomp}, excluding only the boundary
local-time piece $K_{\varepsilon}^{\mathrm{bdy}}$ of~\eqref{eq:Kbdy}.

\begin{theorem}[Interior first-chaos
\texorpdfstring{$H^{-1}$}{H\textasciicircum-1} convergence]
\label{thm:interior}
Let $f\in\mathscr S(\R)$. Then both
$K_{\varepsilon}^{\mathrm{int}}$ and
$K_{\varepsilon}^{\mathrm{int,exact}}=K_{\varepsilon}^{\mathrm{int}}+R_{\varepsilon}^{\mathrm{layer}}$
converge to the same limit
\[
 K_{\varepsilon}^{\mathrm{int}},\,K_{\varepsilon}^{\mathrm{int,exact}}
 \;\xrightarrow[\varepsilon\to 0]{}\;K^{\mathrm{int}}
 \quad\text{in }L^{2}\bigl(\Omega_{Y};\,H^{-1}(\R,dw)\bigr),
\]
with
\begin{equation}\label{eq:Kint-limit}
 K^{\mathrm{int}}(w\mid Y)
 =\tfrac12\bigl[T_{1}^{*}(w\mid Y)+B(w)L(w)-\partial_{w}\bigl(C(w)L(w)\bigr)\bigr],
\end{equation}
where $T_{1}^{*}(w\mid Y)=\mathbf 1_{w\ge 0}\int_{w}^{M_{t}}A(r)L(r)\,dr$,
$L=L_{t}^{\cdot}(|Y|)$, $M_{t}=\sup_{[0,t]}|Y_{s}|$, and $A,B,C$ are
as in~\eqref{eq:ABC-limit}.
\end{theorem}

\begin{remark}[Relation to the
\texorpdfstring{$L^{2}(\Omega)$}{L\textasciicircum 2(Omega)} obstruction]
\label{rem:relation-to-obstruction}
\Cref{thm:interior} does not use the nondegeneracy
hypothesis~\eqref{eq:nondegen}; the convergence holds for every
$f\in\mathscr S(\R)$. The role of~\eqref{eq:nondegen} is to certify
that the present $H^{-1}$-convergent sector is the same sector
whose $L^{2}(\Omega)$-variance diverges in~\cite{Fontes-EnergyFirst}:
under~\eqref{eq:nondegen}, $C\not\equiv 0$ on $(0,\infty)$, so the
$T_{3}^{*}=-\partial_{w}(CL)$ summand in~\eqref{eq:Kint-limit} is
non-trivial, and the divergence in~\eqref{eq:obstruction} is
precisely the $L^{2}(\R)$-norm failure of the same summand. Without
\eqref{eq:nondegen}, the limit~\eqref{eq:Kint-limit} still exists in
$H^{-1}$; it is simply not witnessed by an $L^{2}$ obstruction.
\end{remark}

\begin{proof}
By the triangle inequality, $K_{\varepsilon}^{\mathrm{int,exact}}-K^{\mathrm{int}}
=(K_{\varepsilon}^{\mathrm{int}}-K^{\mathrm{int}})+R_{\varepsilon}^{\mathrm{layer}}$,
and $R_{\varepsilon}^{\mathrm{layer}}\to 0$ in
$L^{2}(\Omega_{Y};L^{2}(\R))\hookrightarrow L^{2}(\Omega_{Y};H^{-1}(\R))$
by \cref{lem:bdy-layer-vanish}. It therefore suffices to prove
$K_{\varepsilon}^{\mathrm{int}}\to K^{\mathrm{int}}$ in
$L^{2}(\Omega_{Y};H^{-1}(\R))$. Throughout we use the continuous
inclusion $L^{2}(\R)\hookrightarrow H^{-1}(\R)$
($\|\cdot\|_{H^{-1}}\le\|\cdot\|_{L^{2}}$).

Decompose
$K_{\varepsilon}^{\mathrm{int}}-K^{\mathrm{int}}
=\tfrac12\sum_{j=1}^{3}\Delta_{j}^{\varepsilon}$
with $\Delta_{j}^{\varepsilon}=T_{j}^{\varepsilon}-T_{j}^{*}$.

\smallskip\noindent\emph{(a) $T_{1}^{\varepsilon}\to T_{1}^{*}$ in
$L^{2}(\Omega_{Y};L^{2}(\R))$.}
Write $\Delta_{1}^{\varepsilon}(w)=I_{\varepsilon}(w)+J_{\varepsilon}(w)$ with
\[
 I_{\varepsilon}(w)=\!\int_{0}^{\infty}\![A_{\varepsilon}(r)-A(r)]L(r)\,\xi_{r,\varepsilon}(w)\,dr,
 \quad
 J_{\varepsilon}(w)=\!\int_{0}^{\infty}\!A(r)L(r)\,[\xi_{r,\varepsilon}(w)-\sigma_{r}(w)]\,dr.
\]
Since $L$ has compact support in $[0,M_{t}]$, the inner integrals
are over $r\in[0,M_{t}]$. By heat-semigroup time-continuity
(\cref{lem:heat-rate}), $A_{\varepsilon}(r)\to A(r)$ uniformly in
$r\in[0,M_{t}]$ a.s.\ in $Y$, while the bulk bound
$\|\xi_{r,\varepsilon}\|_{L^{2}(\R)}^{2}\le|r|+C_{\rho}\varepsilon$
of~\cite{Fontes-EnergyFirst} gives
$\|\xi_{r,\varepsilon}\|_{L^{2}}\le(M_{t}+C_{\rho})^{1/2}$.
Minkowski in $w$ then yields
$\|I_{\varepsilon}\|_{L^{2}(\R)}\le\|A_{\varepsilon}-A\|_{L^{\infty}([0,M_{t}])}
\cdot t\cdot(M_{t}+C_{\rho})^{1/2}\to 0$ a.s.\ in $Y$.
For $J_{\varepsilon}$, the scenery-derivative kernel rate
$\|\xi_{r,\varepsilon}-\sigma_{r}\|_{L^{2}(\R)}\le C'_{\rho}\,\varepsilon^{1/2}$
of~\cite{Fontes-EnergyFirst} gives
$\|J_{\varepsilon}\|_{L^{2}(\R)}\le\|A\|_{\infty}\cdot t\cdot C'_{\rho}\,\varepsilon^{1/2}\to 0$
a.s.\ in $Y$. Both bounds are dominated by $C(f,\rho)\,t\cdot(1+M_{t}^{1/2})$
which is in $L^{2}(\Omega_{Y})$ since $\E_{Y}[M_{t}]<\infty$
classically; dominated convergence concludes
$\E_{Y}\|\Delta_{1}^{\varepsilon}\|_{L^{2}(\R)}^{2}\to 0$.

\smallskip\noindent\emph{(b) $T_{2}^{\varepsilon}\to BL$ in
$L^{2}(\Omega_{Y};L^{2}(\R))$.}
$\Delta_{2}^{\varepsilon}=\rho_{\varepsilon}*((B_{\varepsilon}-B)L)
+(\rho_{\varepsilon}*BL-BL)$. By Young's inequality and
$\|B_{\varepsilon}-B\|_{L^{\infty}([0,M_{t}])}\to 0$
(\cref{lem:heat-rate}),
$\|\rho_{\varepsilon}*((B_{\varepsilon}-B)L)\|_{L^{2}}\le
\|B_{\varepsilon}-B\|_{\infty}\|L\|_{L^{2}}\to 0$ a.s.\ in $Y$. By
standard $L^{2}$-mollifier convergence,
$\|\rho_{\varepsilon}*BL-BL\|_{L^{2}}\to 0$ a.s.\ in $Y$. Both terms
are dominated by $\|B\|_{\infty}\|L\|_{L^{2}}\in L^{2}(\Omega_{Y})$;
dominated convergence concludes.

\smallskip\noindent\emph{(c) $T_{3}^{\varepsilon}\to-\partial_{w}(CL)$
in $L^{2}(\Omega_{Y};H^{-1}(\R))$.}
The operator $\partial_{w}\colon L^{2}(\R)\to H^{-1}(\R)$ has norm
$\le 1$, since for $\phi\in H^{1}(\R)$,
$|\langle\partial_{w}f,\phi\rangle|=|\langle f,\phi'\rangle|\le\|f\|_{L^{2}}\|\phi'\|_{L^{2}}\le\|f\|_{L^{2}}\|\phi\|_{H^{1}}$.
Hence
$\|\Delta_{3}^{\varepsilon}\|_{H^{-1}}=\|\partial_{w}(\rho_{\varepsilon}*C_{\varepsilon}L-CL)\|_{H^{-1}}
\le\|\rho_{\varepsilon}*C_{\varepsilon}L-CL\|_{L^{2}}\to 0$ a.s.\ in $Y$
by the same Young + mollifier argument as (b) applied to $g=CL$.
Domination by $\|C\|_{\infty}\|L\|_{L^{2}}\in L^{2}(\Omega_{Y})$
finishes via dominated convergence.

\smallskip Combining (a)--(c) and applying
$\|\cdot\|_{H^{-1}}\le\|\cdot\|_{L^{2}}$ to (a) and (b),
$\|K_{\varepsilon}^{\mathrm{int}}-K^{\mathrm{int}}\|_{L^{2}(\Omega_{Y};H^{-1}(\R))}\to 0$.
\end{proof}

\subsection*{Quantitative rate (conditional).}

\begin{proposition}[Quantitative interior rate]\label{prop:interior-rate}
Assume \cref{assu:lt-moment}. Then for every
$\eta\in(0,1/2)$,
\begin{equation}\label{eq:int-rate}
 \bigl\|K_{\varepsilon}^{\mathrm{int,exact}}-K^{\mathrm{int}}\bigr\|_{L^{2}(\Omega_{Y};H^{-1}(\R))}
 \;\le\;C(f,t,\rho,\eta)\,\varepsilon^{1/2-\eta}.
\end{equation}
The same bound holds with $K_{\varepsilon}^{\mathrm{int}}$ in place
of $K_{\varepsilon}^{\mathrm{int,exact}}$.
\end{proposition}

\begin{proof}
By \cref{thm:interior},
$K_{\varepsilon}^{\mathrm{int,exact}}-K^{\mathrm{int}}
=\tfrac12\sum_{j=1}^{3}\Delta_{j}^{\varepsilon}+R_{\varepsilon}^{\mathrm{layer}}$.
The residual is controlled by \cref{lem:bdy-layer-vanish}
(eq.~\eqref{eq:R-layer-bound}):
$\|R_{\varepsilon}^{\mathrm{layer}}\|_{L^{2}(\Omega_{Y};L^{2}(\R))}^{2}
\le C(f,\rho)^{2}\,\varepsilon\,\E_{Y}[\|L\|_{\infty}^{2}]$,
so $\|R_{\varepsilon}^{\mathrm{layer}}\|_{L^{2}(\Omega_{Y};H^{-1}(\R))}\le
C(f,\rho)\,\varepsilon^{1/2}\cdot\E_{Y}[\|L\|_{\infty}^{2}]^{1/2}$.
We refine the three $\Delta_{j}^{\varepsilon}$-estimates of
\cref{thm:interior}.

\smallskip\noindent\emph{(a) Rate for
$\Delta_{1}^{\varepsilon}=T_{1}^{\varepsilon}-T_{1}^{*}$.}
Continuing the splitting
$\Delta_{1}^{\varepsilon}=I_{\varepsilon}+J_{\varepsilon}$ of
\cref{thm:interior}(a): by \cref{lem:heat-rate},
$\|A_{\varepsilon}-A\|_{L^{\infty}([0,M_{t}])}\le C_{f}\,\varepsilon$
(uniform); by \cite{Fontes-EnergyFirst} (scenery-derivative kernel
rate), $\|\xi_{r,\varepsilon}-\sigma_{r}\|_{L^{2}(\R)}\le
C'_{\rho}\,\varepsilon^{1/2}$ uniformly in $r$. Hence
$\|I_{\varepsilon}\|_{L^{2}(\R)}\le C_{f}\,\varepsilon\cdot t\cdot(M_{t}+C_{\rho})^{1/2}$
and $\|J_{\varepsilon}\|_{L^{2}(\R)}\le\|A\|_{\infty}\cdot t\cdot C'_{\rho}\,\varepsilon^{1/2}$,
a.s.\ in $Y$. Combining and squaring,
$\|\Delta_{1}^{\varepsilon}\|_{L^{2}(\R)}^{2}\le
C(f,\rho)^{2}\,(t^{2}(1+M_{t}))\,\varepsilon$, so
$\E_{Y}\|\Delta_{1}^{\varepsilon}\|_{L^{2}(\R)}^{2}\le
C(f,\rho)^{2}\,t^{2}(1+\E_{Y}[M_{t}])\,\varepsilon$ (classical
$\E_{Y}[M_{t}]<\infty$).

\smallskip\noindent\emph{(b) Rate for
$\Delta_{2}^{\varepsilon}=T_{2}^{\varepsilon}-BL$.}
$\Delta_{2}^{\varepsilon}=\rho_{\varepsilon}*((B_{\varepsilon}-B)L)
+(\rho_{\varepsilon}*BL-BL)$. By Young,
$\|\rho_{\varepsilon}*((B_{\varepsilon}-B)L)\|_{L^{2}}\le
\|B_{\varepsilon}-B\|_{L^{\infty}([0,M_{t}])}\,\|L\|_{L^{2}}\le C_{f}\,\varepsilon\,\|L\|_{L^{2}}$.
For the second summand, the standard mollifier estimate against
$H^{s}$ for $s\in(0,1]$ states
$\|\rho_{\varepsilon}*g-g\|_{L^{2}}\le C_{\rho}\,\varepsilon^{s}\,\|g\|_{H^{s}}$
(Plancherel: $\int(1+k^{2})^{-s}(1+k^{2})^{s}|\hat g(k)|^{2}|\hat\rho(\varepsilon k)-1|^{2}dk$
with $|\hat\rho(\varepsilon k)-1|\le c_{\rho}\,\varepsilon^{s}(1+k^{2})^{s/2}$
for $s\le 1$). Apply with $g=BL$ and $s=1/2-\eta$: by
\cref{lem:heat-rate}, $B\in C_{b}^{\infty}([0,\infty))$, so
$\|BL\|_{H^{1/2-\eta}}\le C_{f}\,\|L\|_{H^{1/2-\eta}}$, and
\cref{assu:lt-moment} gives
$\E_{Y}\|L\|_{H^{1/2-\eta}}^{2}<\infty$. Combining,
$\E_{Y}\|\Delta_{2}^{\varepsilon}\|_{L^{2}(\R)}^{2}\le
C(f,\rho)^{2}\,\varepsilon^{1-2\eta}$ for every $\eta\in(0,1/2)$.

\smallskip\noindent\emph{(c) Rate for
$\Delta_{3}^{\varepsilon}=T_{3}^{\varepsilon}-T_{3}^{*}$ in
$H^{-1}$.}
Set $u_{\varepsilon}=\rho_{\varepsilon}*(C_{\varepsilon}L)$. By the
same argument as (b) applied to $g=CL$ and $s=1/2-\eta$,
$\E_{Y}\|u_{\varepsilon}-CL\|_{L^{2}(\R)}^{2}\le
C(f,\rho)^{2}\,\varepsilon^{1-2\eta}$. Continuity of
$\partial_{w}\colon L^{2}(\R)\to H^{-1}(\R)$ with norm $\le 1$ gives
$\E_{Y}\|\Delta_{3}^{\varepsilon}\|_{H^{-1}(\R)}^{2}\le
\E_{Y}\|u_{\varepsilon}-CL\|_{L^{2}(\R)}^{2}\le
C(f,\rho)^{2}\,\varepsilon^{1-2\eta}$.

\smallskip Combining
$\|\cdot\|_{H^{-1}}\le\|\cdot\|_{L^{2}}$ on (a) and (b),
$\|\Delta_{j}^{\varepsilon}\|_{L^{2}(\Omega_{Y};H^{-1}(\R))}\le
C(f,t,\rho)\,\varepsilon^{1/2-\eta}$ for each $j=1,2,3$, and the
$R_{\varepsilon}^{\mathrm{layer}}$ contribution above is
$O(\varepsilon^{1/2})\le O(\varepsilon^{1/2-\eta})$. Summing
yields~\eqref{eq:int-rate}.
\end{proof}

\begin{remark}[Why the rate is sharp at
\texorpdfstring{$\varepsilon^{1/2}$}{epsilon\textasciicircum 1/2}]
The Bouleau--Yor regularity $L\in H^{1/2-\eta}$ is sharp at
$s=1/2$, and the mollifier rate
$\|\rho_{\varepsilon}*g-g\|_{L^{2}}\sim\varepsilon^{s}$ for
$g\in H^{s}$ is sharp at the regularity index of $g$. The
$\eta$-loss is intrinsic to the Brownian local-time regularity, not
an artefact of the proof.
\end{remark}

%==============================================================
\section{Boundary first-chaos convergence}\label{sec:boundary}

\begin{theorem}[Boundary first-chaos
\texorpdfstring{$H^{-1/2-\delta}$}{H\textasciicircum-1/2-delta} convergence]
\label{thm:boundary}
For every $f\in\mathscr S(\R)$ and every $\delta>0$,
\[
 K_{\varepsilon}^{\mathrm{bdy}}\;\xrightarrow[\varepsilon\to 0]{}\;L_{t}^{0}(Y)\cdot f'(x)\,\delta_{0}
 \quad\text{in }L^{2}\bigl(\Omega_{Y};\,H^{-1/2-\delta}(\R)\bigr),
\]
with rate
\[
 \bigl\|K_{\varepsilon}^{\mathrm{bdy}}-L_{t}^{0}(Y)f'(x)\delta_{0}\bigr\|_{L^{2}(\Omega_{Y};H^{-1/2-\delta}(\R))}
 \le C_{f,\rho}\,\varepsilon^{\delta\wedge 1/2}.
\]
No $\delta'_{0}$ contribution arises (\cref{lem:bdy-cov-zero},
\cref{rem:bdy-coeff}).
\end{theorem}

\begin{proof}
By~\eqref{eq:Kbdy}, $K_{\varepsilon}^{\mathrm{bdy}}=L_{t}^{0}(Y)\,C_{\varepsilon}(0)\,\rho_{\varepsilon}$.
Split
\begin{equation}\label{eq:bdy-split}
 K_{\varepsilon}^{\mathrm{bdy}}-L_{t}^{0}(Y)\,f'(x)\,\delta_{0}
 \;=\;L_{t}^{0}(Y)\,\bigl(C_{\varepsilon}(0)-f'(x)\bigr)\,\rho_{\varepsilon}
 \;+\;L_{t}^{0}(Y)\,f'(x)\,\bigl(\rho_{\varepsilon}-\delta_{0}\bigr).
\end{equation}

\emph{Step 1: coefficient rate.}
$C_{\varepsilon}(0)=(P_{\sigma_{\varepsilon}^{2}(0)}f')(x)$, and
$\sigma_{\varepsilon}^{2}(0)=I_{1}(0)-\tfrac12 I_{2}(0)$ in the
notation of \cref{lem:bdy-cov-zero}; from
$I_{1}(0)=\int\rho_{\varepsilon}(z)|z|\,dz=c_{\rho}\varepsilon$ and
$I_{2}(0)=O(\varepsilon)$, $\sigma_{\varepsilon}^{2}(0)=O(\varepsilon)$.
By \cref{lem:heat-rate}, $|C_{\varepsilon}(0)-f'(x)|\le
\tfrac12\sigma_{\varepsilon}^{2}(0)\,\|f'''\|_{\infty}\le C_{f}\,\varepsilon$.

\emph{Step 2: distribution rate.}
For $s>1/2$, Plancherel and $\hat{\rho}_{\varepsilon}(k)=\hat\rho(\varepsilon k)$:
\[
 \|\rho_{\varepsilon}-\delta_{0}\|_{H^{-s}}^{2}
 \;=\;\int_{\R}(1+k^{2})^{-s}\,|\hat\rho(\varepsilon k)-1|^{2}\,dk.
\]
Using $|\hat\rho(\varepsilon k)-1|\le\min(2,c_{\rho}\,\varepsilon^{2}k^{2})$
(the bound $c_{\rho}\,\varepsilon^{2}k^{2}$ holds for $|\varepsilon k|\le 1$
by Taylor with $\hat\rho(0)=1$, $\hat\rho'(0)=0$ by evenness; the
bound $2$ holds always), split the integral at $|k|=1/\varepsilon$:
\[
 \|\rho_{\varepsilon}-\delta_{0}\|_{H^{-s}}^{2}
 \;\le\;c_{\rho}^{2}\,\varepsilon^{4}\!\int_{|k|\le 1/\varepsilon}\!(1+k^{2})^{-s}k^{4}\,dk
 \;+\;4\!\int_{|k|>1/\varepsilon}\!(1+k^{2})^{-s}\,dk.
\]
For $1/2<s<5/2$, the first integral is $O(\varepsilon^{-(5-2s)})$ so
the first term is $O(\varepsilon^{4-(5-2s)})=O(\varepsilon^{2s-1})$;
the second is $O(\varepsilon^{2s-1})$. Hence
$\|\rho_{\varepsilon}-\delta_{0}\|_{H^{-s}}\le C\,\varepsilon^{s-1/2}$
for $s\in(1/2,5/2)$. Applying with $s=1/2+\delta$,
$\|\rho_{\varepsilon}-\delta_{0}\|_{H^{-1/2-\delta}}\le C\,\varepsilon^{\delta}$.

\emph{Step 3: assemble.}
Since $\|\rho_{\varepsilon}\|_{H^{-1/2-\delta}}\le
\|\rho_{\varepsilon}\|_{H^{-1/2-\delta}}-\|\delta_{0}\|_{H^{-1/2-\delta}}+\|\delta_{0}\|_{H^{-1/2-\delta}}\le C(\delta)$
(uniform in $\varepsilon$ for small $\varepsilon$, with
$\|\delta_{0}\|_{H^{-1/2-\delta}}<\infty$ since $\delta_{0}\in H^{-1/2-\delta}$),
take $L^{2}(\Omega_{Y};H^{-1/2-\delta}(\R))$-norm
in~\eqref{eq:bdy-split} via the triangle inequality and the
random-scalar pull-out:
\begin{align*}
 \bigl\|K_{\varepsilon}^{\mathrm{bdy}}-L_{t}^{0}(Y)f'(x)\delta_{0}\bigr\|_{L^{2}(\Omega_{Y};H^{-1/2-\delta})}
 &\le\|L_{t}^{0}(Y)\|_{L^{2}(\Omega_{Y})}\,|C_{\varepsilon}(0)-f'(x)|\,\|\rho_{\varepsilon}\|_{H^{-1/2-\delta}}\\
 &\quad+\|L_{t}^{0}(Y)\|_{L^{2}(\Omega_{Y})}\,|f'(x)|\,\|\rho_{\varepsilon}-\delta_{0}\|_{H^{-1/2-\delta}}\\
 &\le C_{\delta}\,C_{f}\,\varepsilon\,\|L_{t}^{0}(Y)\|_{L^{2}(\Omega_{Y})}
  + C\,|f'(x)|\,\varepsilon^{\delta}\,\|L_{t}^{0}(Y)\|_{L^{2}(\Omega_{Y})}.
\end{align*}
$\E[L_{t}^{0}(Y)^{2}]<\infty$ is classical (Brownian local-time
second moment at zero,
\cite[Theorem 6.1.1]{MarcusRosen}). Capping the two
contributions $\varepsilon$ and $\varepsilon^{\delta}$ by their
slower component $\varepsilon^{\delta\wedge 1/2}$ (conservative
rate, accommodating later $H^{-1/2-\delta}\hookrightarrow H^{-1}$
use in \cref{cor:combined}, valid for $\delta\in(0,1/2)$) yields
the claimed bound.
\end{proof}

\begin{remark}[Why
\texorpdfstring{$H^{-1/2-\delta}$}{H\textasciicircum-1/2-delta},
not \texorpdfstring{$H^{-1}$}{H\textasciicircum-1}, for the boundary]
$\delta_{0}\in H^{s}(\R)$ iff $s<-1/2$, so $\delta_{0}\notin H^{-1/2}$
and $\delta_{0}\in H^{-1/2-\delta}$ for every $\delta>0$. Asking
for convergence in $H^{-1/2}$ is impossible at the boundary because
the limit object is not in that space. The natural boundary space is
$H^{-1/2-\delta}$. Since $H^{-1/2-\delta}\subset H^{-1}$ for
$\delta\in(0,1/2)$, the boundary kernel also converges in $H^{-1}$,
which is used in~\cref{cor:combined}.
\end{remark}

%==============================================================
\section{Combined first-chaos dual closure}\label{sec:combined}

\begin{corollary}[Combined first-chaos
\texorpdfstring{$H^{-1}$}{H\textasciicircum-1} convergence]
\label{cor:combined}
For every $f\in\mathscr S(\R)$ and every $\delta\in(0,1/2)$, the
full first-$X$-chaos Malliavin kernel
\[
 K_{\varepsilon}
 \;=\;K_{\varepsilon}^{\mathrm{bdy}}+K_{\varepsilon}^{\mathrm{int,exact}}
 \;=\;K_{\varepsilon}^{\mathrm{bdy}}+K_{\varepsilon}^{\mathrm{int}}+R_{\varepsilon}^{\mathrm{layer}}
\]
satisfies
\[
 K_{\varepsilon}\;\xrightarrow[\varepsilon\to 0]{}\;K
 \;\coloneqq\;L_{t}^{0}(Y)f'(x)\delta_{0}+K^{\mathrm{int}}
 \quad\text{in }L^{2}\bigl(\Omega_{Y};\,H^{-1}(\R,dw)\bigr).
\]
The boundary contribution converges in $H^{-1/2-\delta}\subset H^{-1}$
(continuous embedding for $\delta\in(0,1/2)$); the interior
contribution $K_{\varepsilon}^{\mathrm{int,exact}}$ converges in
$L^{2}(\R)\subset H^{-1}$ via \cref{thm:interior,lem:bdy-layer-vanish}.
Under \cref{assu:lt-moment}, the combined rate is
$\varepsilon^{(1/2-\eta)\wedge\delta}$ for every $\eta\in(0,1/2)$
and $\delta\in(0,1/2)$; choosing $\delta=1/2-\eta$ optimises to
$\varepsilon^{1/2-\eta}$.
\end{corollary}

\begin{proof}
Combine \cref{thm:interior,thm:boundary}. The embedding
$H^{-1/2-\delta}\hookrightarrow H^{-1}$ for $\delta\in(0,1/2)$ is
the standard continuous Sobolev inclusion
($H^{s}\subset H^{s'}$ for $s\ge s'$, applied with
$s=-1/2-\delta\ge -1=s'$ iff $\delta\le 1/2$); the operator norm of
the embedding is $\le 1$. The interior $L^{2}(\R)\subset H^{-1}(\R)$
embedding is even more direct ($\|\cdot\|_{H^{-1}}\le\|\cdot\|_{L^{2}}$).
\end{proof}

\begin{corollary}[First-chaos dual closure of
\texorpdfstring{$\mathrm{Proj}_{1}^{X}C_{\IBM,\varepsilon}$}{Proj\textasciicircum X\_1 C\_\{IBM,eps\}}]
\label{cor:proj1-closure}
Let $\Energy_{\IBM,1,*}^{(4),-}=H^{-1}(\R,dw)$. The first-scenery-chaos
projection $\mathrm{Proj}_{1}^{X}C_{\IBM,\varepsilon}(f;t)$, viewed
as a $Y$-conditioned first-chaos element with conditional Malliavin
kernel in $L^{2}(\Omega_{Y};\Energy_{\IBM,1,*}^{(4),-})$, converges
to an explicit $Y$-conditioned first-chaos element with kernel $K$
given by \cref{cor:combined} as $\varepsilon\to 0$. Pairing
against any deterministic $\phi\in H^{1}(\R)$ gives
\[
 \bigl\langle\mathrm{Proj}_{1}^{X}C_{\IBM,\varepsilon}(f;t),\phi\bigr\rangle_{1,*}
 \;\xrightarrow[\varepsilon\to 0]{}\;
 \E_{Y}\!\bigl[\langle K(\cdot\mid Y),\phi\rangle_{H^{-1},H^{1}}\bigr].
\]
Under \cref{assu:lt-moment}, the rate is
$\varepsilon^{(1/2-\eta)\wedge\delta}$ for every $\eta\in(0,1/2)$
and $\delta\in(0,1/2)$; choosing $\delta=1/2-\eta$ optimises to
$\varepsilon^{1/2-\eta}$.
\end{corollary}

\subsection*{Scope: what this closes and what it does not.}\label{sec:scope}
\Cref{cor:combined,cor:proj1-closure} close the
$L^{2}(\Omega)$-divergent first scenery-chaos sector of
$C_{\IBM,\varepsilon}(f;t)$ in the first-$X$-chaos block of the
quartic dual. They do \emph{not} prove:
\begin{itemize}
\item that $C_{\IBM}(f;t)$ exists in the full quartic dual
$\Energy_{\IBM}^{(4),-}$: the $(m,n)$-bidegree projections with
$m\ge 2$, including the second-energy-contraction sector
controlled by the contraction-renormalised correction candidate
of~\cite{Fontes-EnergyFirst} (whose Gaussian-coordinate realisation
is the Wick-square subtraction), are not addressed;
\item that the full unmollified IBM It\^o
identity~\eqref{eq:flagship} survives the $\varepsilon\to 0$ limit
in any dual topology;
\item that the Brownian-time semigroup correction
$u(t,x)-f(x)$ of~\cite{Fontes-EnergyFirst} arises from the
first-chaos sector: it cannot, because first chaos has mean zero, so
$\E[\mathrm{Proj}_{1}^{X}C_{\IBM}(f;t)]=0$.
\end{itemize}

Filtration-level obstructions (Tsirelson-type non-Brownian noises
etc.) are an independent direction and play no role in any argument
above.

%==============================================================
\section{Lemmas imported from EnergyFirst}\label{sec:imports}

We cite (do not re-prove) the following inputs from
\cite{Fontes-EnergyFirst}, referring to them by descriptive name
rather than internal label number:

\begin{enumerate}
\item \emph{Mollifier $L^{2}$-rate.}
$\E[(X^{\varepsilon}_{x}-X_{x})^{2}]\le C_{\rho}\varepsilon$
uniformly in $x\in\R$. Used implicitly in $C_{\varepsilon}\to C$ and
$\sigma_{\varepsilon}^{2}(0)=O(\varepsilon)$.
\item \emph{Scenery-derivative kernel $L^{2}$-bound and rate.}
$\|H_{\varepsilon}(y,\cdot)\|_{L^{2}}^{2}\le|y|+C_{\rho}\varepsilon$
and $\|H_{\varepsilon}(y,\cdot)-\sigma_{y}\|_{L^{2}}^{2}\le
C'_{\rho}\varepsilon$. Used in \cref{thm:interior}(a) and
\cref{prop:interior-rate}(a).
\item \emph{Pointwise variance of $\partial_{y}X^{\varepsilon}$.}
$\Var(\partial_{y}X^{\varepsilon}(y))=c_{\rho}^{(1)}/\varepsilon$ in
the bulk. Used in \cref{sec:audit} bulk-covariance computation.
\item \emph{Kernel decomposition (\cite{Fontes-EnergyFirst}
Proposition on $L^{2}$-divergence, Steps 1--4).} The decomposition
$K_{\varepsilon}=K_{\varepsilon}^{\mathrm{bdy}}+\tfrac12(T_{1}^{\varepsilon}+T_{2}^{\varepsilon}+T_{3}^{\varepsilon})$
of~\eqref{eq:Keps-decomp} on which this paper builds. Re-derived in
\cref{sec:audit}.
\item \emph{Local-time non-Sobolev regularity.}
$L_{t}^{\cdot}(Y)\notin H^{1}_{\mathrm{loc}}$ a.s.\ on
$\{M_{t}>b\}$. Used as the sharp lower bound matching the
$H^{-1}$ upper bound (cf.\ \cref{rem:why-Hminus1}).
\end{enumerate}

\noindent\textbf{New auxiliary lemmas, stated here (proofs deferred):}

\begin{lemma}[Heat-semigroup time-Lipschitz]\label{lem:heat-rate}
For $h\in\mathscr S(\R)$ and $\eta_{1},\eta_{2}\ge 0$,
$\|P_{\eta_{1}}h-P_{\eta_{2}}h\|_{L^{\infty}(\R)}\le
\tfrac12|\eta_{1}-\eta_{2}|\,\|h''\|_{L^{\infty}}$. In particular
$|C_{\varepsilon}(r)-C(r)|\le C_{f}|\sigma_{\varepsilon}^{2}(r)-|r||\le
C_{f}\varepsilon$ uniformly in $|r|\ge c\varepsilon$, and
$|C_{\varepsilon}(0)-f'(x)|\le C_{f}\varepsilon$.
\end{lemma}

\begin{lemma}[Reflected-BM spatial local-time Sobolev regularity]
\label{lem:lt-H12-eta}
For every $t>0$ and $\eta\in(0,1/2)$,
$L_{t}^{\cdot}(|Y|)\in H^{1/2-\eta}([0,M_{t}])$ a.s. Equivalently,
$L_{t}^{x}(|Y|)=L_{t}^{x}(Y)+L_{t}^{-x}(Y)$ on $(0,\infty)$, and
each summand is the spatial Brownian local time of two-sided BM.
\end{lemma}

\Cref{lem:heat-rate} is classical. \Cref{lem:lt-H12-eta} is standard
from the Marcus--Rosen literature on local-time regularity;
\cref{assu:lt-moment} is the integrated form whose precise reference
will appear in the final version.

%==============================================================
\section{Risk checklist}\label{sec:risks}

The following items concern sharpening and validating the
\emph{quantitative} statements (\cref{prop:interior-rate} and the
combined rate in \cref{cor:combined,cor:proj1-closure}). The
\emph{qualitative} theorems
(\cref{thm:interior,thm:boundary,cor:combined})
are supported by proved lemmas
(\cref{lem:bdy-cov-zero,lem:bdy-layer-vanish}) and do not depend on
the items below; only the rate and the explicit constants do.

\begin{risk}[\textbf{1.} Reflected BM local-time regularity]
\label{rsk:reflected-lt}
\Cref{lem:lt-H12-eta} for $|Y|$ rather than $Y$. The decomposition
$L_{t}^{x}(|Y|)=L_{t}^{x}(Y)+L_{t}^{-x}(Y)$ for $x>0$ is standard,
and each summand is $H^{1/2-\eta}_{\mathrm{loc}}$ a.s.\ on
$(0,\infty)$, so $L_{t}^{\cdot}(|Y|)\in H^{1/2-\eta}_{\mathrm{loc}}((0,\infty))$.
At $x=0$ the formula gives $L_{t}^{0}(|Y|)=2L_{t}^{0}(Y)$ (Tanaka),
and the function $x\mapsto L_{t}^{x}(|Y|)$ is continuous and
$H^{1/2-\eta}$ on $[0,M_{t}]$ a.s. \textbf{Verify} the relevant
Marcus--Rosen reference and confirm uniform constants on
$\{M_{t}\le R\}$ before integrating tails.
\end{risk}

\begin{risk}[\textbf{2.}
\texorpdfstring{$\varepsilon^{1/2-\eta}$}{epsilon\textasciicircum 1/2-eta} rate in
\texorpdfstring{$L^{2}(\Omega_{Y})$}{L\textasciicircum 2(Omega\_Y)}]
\label{rsk:rate-L2Y}
\Cref{prop:interior-rate} requires \cref{assu:lt-moment}, the
integrated form of $L_{t}^{\cdot}(|Y|)\in H^{1/2-\eta}_{\mathrm{loc}}$.
The classical $\E_{Y}[M_{t}^{p}]<\infty$ for all $p$ is one half;
the other half is uniformity of the Marcus--Rosen constant on
$\{M_{t}\le R\}$ with tail control via Gaussian moments of $M_{t}$.
\textbf{Verify}.
\end{risk}

\begin{risk}[\textbf{3.} Boundary kernel coefficient]
\label{rsk:bdy-coefficient}
\Cref{lem:bdy-cov-zero} establishes $\Cov(U_{0},V_{0})=0$, killing
the $\xi_{0,\varepsilon}$ piece in $\E_{X}[D_{w}g_{\varepsilon}'(0)]$
and ruling out a $\delta'_{0}$ contribution. \textbf{Verify} that
the calculation is robust under the choice of two-sided scenery
convention $\Cov(X_{z},X_{z'})=\tfrac12(|z|+|z'|-|z-z'|)$, and that
the integration by parts in the proof respects compact support of
$\rho_{\varepsilon}$ (no boundary terms).
\end{risk}

\begin{risk}[\textbf{4.} ``Quartic'' versus Wiener chaos degree]
\label{rsk:quartic-label}
\Cref{rem:quartic-not-chaos} states that the quartic label refers
to operator-energy / test-space order, not Wiener chaos degree.
\textbf{Audit} every occurrence of ``quartic'' to ensure consistency
with this convention; in particular check the introduction and the
abstract do not encourage the misreading that the first-chaos
sector lies in the fourth Wiener chaos.
\end{risk}

\begin{risk}[\textbf{5.} Bochner measurability into
\texorpdfstring{$H^{-1}$}{H\textasciicircum-1}]
\label{rsk:bochner}
The limit kernel $K\in L^{2}(\Omega_{Y};H^{-1}(\R))$ must be
Bochner-measurable. Each $K_{\varepsilon}$ is Bochner-measurable in
$L^{2}(\Omega_{Y};L^{2}(\R))$; $L^{2}\hookrightarrow H^{-1}$
continuous, so $K_{\varepsilon}\in L^{2}(\Omega_{Y};H^{-1}(\R))$.
The limit is strongly measurable as a Cauchy limit. Standard but
worth recording.
\end{risk}

\begin{risk}[\textbf{6.} Pairing identity for first-chaos elements]
\label{rsk:pairing}
\Cref{def:first-chaos-block} writes the pairing as
$\E_{Y}\langle K(\cdot\mid Y),\phi\rangle_{H^{-1},H^{1}}$.
\textbf{Verify} the identity
$\E[I_{1}^{X}(K)I_{1}^{X}(\phi)\mid Y]=\langle K(\cdot\mid Y),\phi\rangle$
in the form required: at finite $\varepsilon$ this is the
$L^{2}(\R)$-pairing via the conditional Wiener isometry on the
scenery factor; in the limit, $K(\cdot\mid Y)\in H^{-1}(\R)$ and the
pairing extends by duality since $\phi\in H^{1}(\R)$. Confirm no
measurability or domain subtleties.
\end{risk}

\begin{risk}[\textbf{7.} Heat-semigroup boundary at
\texorpdfstring{$r\to 0^{+}$}{r->0+}]
\label{rsk:heat-bdy}
$C(r)=(P_{|r|}f')(x)$ is $C^{\infty}$ at $r>0$ but the limit
$C(0)=f'(x)$ is only continuous in $r$ from above. \textbf{Verify}
that $C,B,A$ extend to $C^{\infty}([0,\infty))$ via the heat
semigroup acting on Schwartz functions, so that $T_{1}^{*},BL,CL$
are genuine elements of $L^{2}([0,M_{t}])$ a.s.\ and that
$\partial_{w}(CL)$ in \cref{thm:interior}(c) is a well-defined
$H^{-1}(\R)$ element.
\end{risk}

\smallskip
\noindent\textbf{Resolution priority.}
\Cref{rsk:reflected-lt,rsk:rate-L2Y} are
load-bearing for the quantitative rate
\cref{prop:interior-rate} and its consequences in
\cref{cor:combined,cor:proj1-closure}; they do not affect the
qualitative convergence
\cref{thm:interior,thm:boundary,cor:combined}, which rests on the
proved \cref{lem:bdy-cov-zero,lem:bdy-layer-vanish} and classical
local-time facts. \Cref{rsk:bdy-coefficient} is resolved subject to
checking the scenery convention. \Cref{rsk:quartic-label} is
editorial. \Cref{rsk:bochner,rsk:pairing,rsk:heat-bdy} are routine.

%==============================================================
\begin{thebibliography}{9}
\bibitem{Fontes-EnergyFirst}
R.~Fontes, \emph{Energy-First Stochastic Calculus: Iterated Brownian
Motion as a Test Case}, manuscript, \today.

\bibitem{MarcusRosen}
M.~B.~Marcus and J.~Rosen, \emph{Markov Processes, Gaussian
Processes, and Local Times}, Cambridge Studies in Advanced
Mathematics \textbf{100}, Cambridge University Press, 2006.
\end{thebibliography}

\end{document}
